% © 2026 Ing. David Jaroš — CC BY-NC-ND 4.0
% Licensed under CC BY-NC-ND 4.0. See LICENSE.md.

\ifdefined\INCLUDEMODE
\else
\documentclass[12pt,a4paper]{article}
\usepackage[a4paper,margin=2.5cm]{geometry}
\usepackage{amsmath,amssymb,amsthm}
\usepackage{mathtools}
\usepackage{xcolor}
\usepackage{mdframed}

\newtheorem{theorem}{Theorem}[section]
\newtheorem{proposition}[theorem]{Proposition}
\newtheorem{lemma}[theorem]{Lemma}
\theoremstyle{remark}
\newtheorem{remark}[theorem]{Remark}

\newmdenv[backgroundcolor=green!15,linecolor=green!60!black,linewidth=1.5pt]{resultbox}
\newmdenv[backgroundcolor=yellow!20,linecolor=orange,linewidth=1.5pt]{gapbox}

\title{GAP-U2a: Derivation of the $A_\psi$ Field Equation}
\author{Ing.\ David Jaroš}
\date{2026-07-14}

\begin{document}
\maketitle
\fi

\section{Inputs from existing repository definitions}

We use:
\begin{itemize}
\item $S_{\mathrm{total}}=\frac1{16\pi G}\int\sqrt{-g}R+S_\Theta$ from
\texttt{papers/UBT\_GR\_Submission.tex},
\item kinetic term $\operatorname{Re}\operatorname{Tr}[(D_\mu\Theta)^\dagger D^\mu\Theta]$,
\item gauge-field kinetic term $-\frac14 F_{\mu\nu}F^{\mu\nu}$ in canonical gauge Lagrangian
\texttt{canonical/interactions/sm\_gauge.tex}.
\end{itemize}

Restrict to one Abelian component $A_\mu$ along the $\psi$-sector with
\[
D_\mu\Theta=\partial_\mu\Theta+i q A_\mu\Theta,
\qquad
F_{\mu\nu}=\partial_\mu A_\nu-\partial_\nu A_\mu.
\]

\section{Euler--Lagrange variation with respect to $A_\mu$}

Consider
\[
\mathcal L = \sqrt{-g}\left[\operatorname{Re}\operatorname{Tr}\big((D_\mu\Theta)^\dagger D^\mu\Theta\big)-\frac14F_{\mu\nu}F^{\mu\nu}\right].
\]
Varying $A_\mu\mapsto A_\mu+\delta A_\mu$ gives
\[
\delta D_\mu\Theta=i q\,\delta A_\mu\Theta,
\qquad
\delta F_{\mu\nu}=\nabla_\mu\delta A_\nu-\nabla_\nu\delta A_\mu.
\]
Hence
\[
\delta S=\int d^4x\sqrt{-g}\left(J^\mu\delta A_\mu-(\nabla_\nu F^{\nu\mu})\delta A_\mu\right)
\]
with
\[
J^\mu = i q\,\operatorname{Re}\operatorname{Tr}\!\left(\Theta^\dagger D^\mu\Theta-(D^\mu\Theta)^\dagger\Theta\right).
\]
Arbitrariness of $\delta A_\mu$ implies
\[
\boxed{\nabla_\nu F^{\nu\mu}=J^\mu.}
\]

\begin{theorem}[Derived connection equation]
From the stated action terms, the $A_\mu$ Euler--Lagrange equation is Maxwell-type:
\[
\nabla_\nu F^{\nu\mu}=J^\mu.
\]
This is outcome \textbf{B} in the GAP-U2a problem statement:
\[
\Delta_g A_\psi + F(A_\psi,\Theta)=0
\]
with source/nonlinear term encoded by $J^\mu$.
\textbf{Verdict: PROVED.}
\end{theorem}

\section{Radial $A_\psi$ reduction}

Assume static spherically symmetric sector,
\[
A_\mu dx^\mu=A_\psi(r)\,d\psi,
\qquad g_{ij}=\Psi^4\delta_{ij},
\]
and source-free channel $J^\psi=0$.
Then
\[
\nabla_i\nabla^i A_\psi=0
\quad\Longleftrightarrow\quad
\partial_r\big(r^2\Psi^2 A_\psi'(r)\big)=0.
\]

\begin{proposition}[When harmonic equation follows]
The equation $\Delta_g A_\psi=0$ follows only after imposing $J^\psi=0$ in addition to the derived Euler--Lagrange equation.
\textbf{Verdict: CONDITIONAL.}
\end{proposition}

\section{Insufficient-action branch}

If one keeps only $\operatorname{Re}\operatorname{Tr}[(D\Theta)^\dagger D\Theta]$ and omits gauge kinetic
$F^2$, variation in $A_\mu$ yields an algebraic current constraint ($J^\mu=0$) and no second-order PDE for $A_\mu$.

\begin{proposition}[No-Laplace consequence without $F^2$ term]
Without a gauge kinetic term, deriving $\Delta_g A_\psi=0$ is impossible.
\textbf{Verdict: PROVED.}
\end{proposition}

\begin{gapbox}
This file derives the equation itself. Closure of constants and physical matching (e.g., $C=M$) is handled separately.
\end{gapbox}

\ifdefined\INCLUDEMODE
\else
\end{document}
\fi
